\section{Datos pequeños y aprendizaje bilingüe: convertir la intuición humana en experimentos comprobables}

La clase utiliza el aprendizaje infantil para plantear la cuestión de la eficiencia con datos, pero no equipara las redes neuronales con el cerebro. Los humanos reciben entrada visual, motora, interacción social, metas a largo plazo y alta estructuración; los modelos de lenguaje suelen recibir tokens offline y una función de pérdida unificada. El valor de la investigación reside en construir condiciones de datos controladas, no en demostrar directamente con pocos tokens un "aprendizaje tipo cerebro".

\subsection{¿Por qué estudiar small model y small data?}

Baby Scale compara infant/ChatGPT con brain/neural network, listando aprendizaje continuo, input grounded, metas, estructura jerárquica y sesgo inductivo latente. Esta comparación actúa como generador de hipótesis: si la interactividad, la diversidad semántica o la estructura del hablante en la entrada infantil difieren de la de las páginas web, se puede diseñar una prueba a nivel de dataset; pero no se deben convertir todas las supervisiones humanas a caracteres de transcripción de texto.

\lecturefigure{slide-039.jpg}{Infantes versus ChatGPT, cerebro versus red neuronal: contraste entre dos sistemas de aprendizaje}{CS25 V6 official deck, p.~39}

\lecturefigure{slide-040.jpg}{Diferencias en continuidad, metas, modalidades y estructura entre el aprendizaje humano y el entrenamiento de LLMs}{CS25 V6 official deck, p.~40}

\lecturefigure{slide-041.jpg}{Motivación para estudiar modelos pequeños y datos pequeños: eficiencia, interpretabilidad y experimentos controlados}{CS25 V6 official deck, p.~41}

\teachervoice{00:18:22--00:20:46，el docente utiliza repetidamente might、potential y question para describir las ventajas del aprendizaje humano. Nota de clase: estas diferencias son ejes experimentales, no la conclusión de que "los datos infantiles son naturalmente superiores a los de internet".}

\begin{warningbox}{¿Qué falta al comparar el número de tokens entre infantes y LLMs?}
Los infantes reciben simultáneamente imágenes, sonido, consecuencias de acciones, retroalimentación social e interacción repetida; el corpus transcrita conserva solo una parte. Las funciones objetivo, la inicialización de parámetros, los conocimientos previos y las tareas de evaluación también difieren entre modelos e infantes. Por tanto, el conteo de tokens puede plantear la cuestión de la eficiencia con datos, pero no representa una cantidad total de supervisión justa.
\end{warningbox}

\subsection{Baby Scale: ¿qué puede entrenar la entrada lingüística de un solo niño?}

Baby Scale utiliza speech dirigida al niño (child-directed speech) de una sola familia del corpus CHILDES, entrena modelos de estilo GPT-2 pequeños y compara en tareas de language modeling, sintaxis, semántica y desarrollo lexical. El diseño central no persigue SOTA, sino observar si el \textbf{origen familiar de los datos} y la \textbf{estructura de exposición infantil} generan diferencias medibles una vez fija la escala del modelo.

\lecturefigure{slide-043.jpg}{Pregunta central de Baby Scale: ¿puede un modelo aprender de la entrada de un solo niño?}{CS25 V6 official deck, p.~43}

\lecturefigure{slide-044.jpg}{Configuración experimental: transcripciones familiares, escala de tokens, modelos y tareas de evaluación}{CS25 V6 official deck, p.~44}

La perplexity es la exponencial de la entropía cruzada:
\[
\operatorname{PPL}=\exp\!\left(-\frac{1}{T}\sum_{t=1}^{T}\log p_{\theta}(x_t\mid x_{<t})\right).
\]
Intuitivamente, aproxima el número "efectivo de candidatos" frente a cada token; valores más bajos indican menos sorpresa ante la secuencia de prueba. Sin embargo, PPL está influenciado por el tokenizer, la distribución de prueba y la longitud; no permite comparación directa entre modelos arbitrarios ni equivale a capacidad gramatical, factual o interactiva.

La página de resultados debe leerse en dos pasos: la diapositiva 45 compara el rendimiento relativo de diferentes datos familiares en múltiples tareas, mientras que la 46 revisa si el tamaño del modelo altera la tendencia. Si un conjunto de datos es fuerte en varias tareas, puede deberse a cobertura lexical, estructura interactiva o diversidad del hablante; si ampliar el modelo no elimina las diferencias, indica que las propiedades de los datos no son simplemente un cuello de botella de capacidad.

\lecturefigure{slide-045.jpg}{Resultados multitarea de Baby Scale: diferentes entradas familiares generan perfiles de capacidades distintos}{CS25 V6 official deck, p.~45}

\lecturefigure{slide-046.jpg}{Interacción entre escala del modelo y origen de datos: aumentar parámetros no uniformiza ni elimina las diferencias}{CS25 V6 official deck, p.~46}

\begin{knowledgebox}{Lectura de la figura: comparar primero tareas, luego escalas}
En primer lugar, compare diferentes datos familiares dentro de cada subplot, en lugar de comparar puntuaciones absolutas entre tareas; los métricas y niveles de dificultad varían por tarea. Luego, observe si el ordenamiento se mantiene estable tras cambiar el tamaño del modelo. Un ordenamiento estable apoya la idea de que "las propiedades de los datos afectan continuamente al rendimiento", pero no basta con la correlación para determinar si es la interacción infantil, la cobertura lexical o la calidad de transcripción la causa; se requieren ablaciones controladas posteriores.
\end{knowledgebox}

\subsection{Más allá de la escala: ¿qué propiedades de los datos predicen el rendimiento?}

La clase lista candidatos como interaction, diversidad del hablante, longitud de contexto, cobertura lexical, proporción child/caregiver speech y divergencia de distribución. Esto puede expresarse desde una perspectiva de regresión simple:
\[
y_{f,m}=\beta_0+\beta_1\log D_f+\beta_2\log N_m
+\sum_k\gamma_k z_{f,k}+\varepsilon_{f,m},
\]
donde $y_{f,m}$ es el rendimiento de la tarea para la familia $f$ y escala del modelo $m$, $D_f$ es el número de tokens, $N_m$ son los parámetros y $z_{f,k}$ representa las propiedades de los datos. Esta ecuación no es el único modelo del artículo, sino una ayuda para distinguir efectos de tamaño (size effect) de efectos de composición (composition effect).

\lecturefigure{slide-047.jpg}{Resultados de Scaling: el impacto del dataset y la escala del modelo en tareas distintas no es uniforme}{CS25 V6 official deck, p.~47}

\lecturefigure{slide-048.jpg}{Más allá del tamaño: variables explicativas candidatas como interacción, vocabulario, hablante y contexto}{CS25 V6 official deck, p.~48}

\lecturefigure{slide-049.jpg}{Conclusión general: los datos familiares con mejor rendimiento son más ricos, interactivos y estructuralmente semánticos}{CS25 V6 official deck, p.~49}

\lecturefigure{slide-050.jpg}{Takeaways de Baby Scale: el aprendizaje a escala infantil es viable pero altamente dependiente de la tarea}{CS25 V6 official deck, p.~50}

\teachervoice{00:22:23--00:25:25，el docente enfatiza que "más grande" no es la única explicación: la proporción child/caregiver, la diversidad interactiva y la estructura lexical pueden afectar el rendimiento; además, los datasets óptimos varían según la tarea.}

\begin{importantbox}{Conclusiones que Baby Scale sí y no apoya}
Apoya que: modelos pequeños pueden aprender patrones lingüísticos no aleatorios a partir de transcripciones de un entorno infantil individual; las diferencias en el origen de los datos dejan huellas relevantes para la tarea. No apoya que: la entrada infantil explique por completo el aprendizaje humano del lenguaje; cualquier corpus familiar sea óptimo necesariamente en modelos grandes; o que las propiedades de los datos correlacionadas se hayan demostrado como mecanismos causales.
\end{importantbox}

Al evaluar este tipo de investigación, también hay que evitar "un único puntaje que cubra todas las capacidades". La función de pérdida de language modeling mide predicciones locales; tareas de grammaticality del tipo BLiMP verifican preferencias estructurales; la similitud lexical revisa la representación lexical; y las medidas de desarrollo se interesan por la trayectoria a medida que crece el dato. Un corpus familiar puede dominar en tareas semánticas pero rezagarse en sintácticas; esto no es una contradicción, sino que indica que las propiedades de los datos soportan capacidades distintas. Experimentos posteriores más fuertes deberían realizar emparejamiento o intervención sobre propiedades altamente correlacionadas, por ejemplo re-muestreando diversidad del hablante manteniendo el número de tokens fijo, desordenando turnos de interacción o controlando cobertura lexical, para ver si los perfiles de tareas cambian según lo predicho. Solo así se puede avanzar de "¿qué dataset es mejor" a "¿qué mecanismo de datos está actuando".

El tamaño de la muestra y el entrenamiento repetido también deben ingresar en la incertidumbre. Los experimentos con modelos pequeños suelen verse afectados por random seed, orden del batch y early stopping; si solo se reporta una ejecución, el ordenamiento del dataset podría ser ruido de optimización. Un reporte ideal debería dar medias multi-seed, intervalos de confianza, tamaños de efecto (effect size) y corrección para comparaciones múltiples, además de explicar si existen superposiciones o diferencias de calidad en las transcripciones a nivel familiar. Para las afirmaciones de desarrollo, también se debe comparar la curva de crecimiento del dato en lugar de mirar solo el checkpoint final. Solo así se puede juzgar si un corpus es consistentemente más efectivo o simplemente converge más rápido bajo un presupuesto de entrenamiento limitado.

\subsection{Bilingual BabyLM: ¿es el bilingüismo interferencia, transferencia o problema de escala de datos?}

La investigación bilingüe desglosa aún más "mezclar dos idiomas" en estructura de exposición. Se puede permitir que cada hablante use solo un idioma (speaker-split), que el hablante cambie entre oraciones (sentence-level) o dentro de la misma oración (word-level code-switching). Si todos los casos se llaman bilingual mixture, no se puede determinar a qué es sensible el modelo: límites lingüísticos, pistas del hablante o volumen total de datos.

\lecturefigure{slide-052.jpg}{Tres preguntas de investigación de Bilingual BabyLM: interferencia, estructura de exposición y escala}{CS25 V6 official deck, p.~52}

\lecturefigure{slide-053.jpg}{Condiciones de entrenamiento: monolingüe, bilingüe aleatorio, cambio en nivel speaker/sentence/word}{CS25 V6 official deck, p.~53}

Sea $q_{\mathrm{en}}$ y $q_{\mathrm{es}}$ la distribución de datos en inglés y español, con peso de mezcla $\lambda$; entonces la distribución de entrenamiento es
\[
q_{\lambda}(x)=\lambda q_{\mathrm{en}}(x)+(1-\lambda)q_{\mathrm{es}}(x).
\]
Sin embargo, speaker-split y code-switching no solo cambian $\lambda$; alteran la distribución conjunta de lenguaje/contexto, hablante y límites de oración. Incluso con el mismo total de tokens, las pistas de límite disponibles para el modelo difieren.

\subsection{Resultado uno: no hay interferencia bilingüe universal}

El diseño experimental anterior desglosa la exposición bilingüe en múltiples condiciones; esta sección responde primero a la pregunta más básica: ¿aprender simultáneamente dos idiomas destruye sistemáticamente las capacidades monolingües? Al leer las diapositivas 55--57, primero se debe fijar el total de datos y el baseline, luego verificar por separado language modeling, sintaxis y semántica; solo si múltiples métricas degeneran significativamente en la misma dirección se puede apoyar una interferencia amplia.

La tabla de PPL en la slide 55 muestra que los modelos multilingües pueden formar distribuciones efectivas en ambos idiomas; las slides 56--57 luego usan grammaticality y word similarity para verificar si solo recuerdan estadísticas locales. Al leer la tabla, primero busque el baseline monolingüe, luego revise los intervalos de confianza de las condiciones bilingües, no solo compare una cifra decimal.

\lecturefigure{slide-055.jpg}{Perplexity de modelos bilingües en inglés y español: pueden modelar ambos idiomas simultáneamente}{CS25 V6 official deck, p.~55}

\lecturefigure{slide-056.jpg}{Grammaticality: comparación de capacidad estructural entre baseline en inglés y condiciones multilingües}{CS25 V6 official deck, p.~56}

\lecturefigure{slide-057.jpg}{Word similarity: contraste de representación semántica entre entrenamiento monolingüe y multilingüe}{CS25 V6 official deck, p.~57}

\begin{knowledgebox}{Lectura de la figura: cómo juzgar interferencia}
La interferencia no debe definirse solo como una condición ligeramente inferior al máximo puntaje. Se deben revisar simultáneamente el baseline monolingüe, el volumen total de datos, los intervalos de confianza y múltiples tareas. Si un modelo bilingüe baja ligeramente en inglés pero aprende español, el tradeoff general difiere de la "interferencia catastrófica"; si tras fijar el total de datos aún hay una caída significativa, esto se acerca más a problemas de capacidad competitiva u optimización.
\end{knowledgebox}

\subsection{Resultado dos: la estructura de exposición tiene un impacto menor al esperado}

Las curvas de grammaticality y word similarity para diferentes condiciones multilingües están globalmente cercanas, indicando que los modelos no dependen fuertemente de pistas humanas comunes como "un hablante por idioma". Este resultado debe interpretarse así: bajo el modelo actual, volumen de datos y métricas, el impacto marginal de la estructura de exposición es limitado; no se puede concluir que code-switching sea irrelevante en todos los idiomas, escalas o tareas generativas.

\lecturefigure{slide-059.jpg}{Rendimiento de grammaticality bajo diferentes estructuras de exposición bilingüe}{CS25 V6 official deck, p.~59}

\lecturefigure{slide-060.jpg}{Rendimiento de word similarity bajo diferentes estructuras de exposición bilingüe}{CS25 V6 official deck, p.~60}

\subsection{Resultado tres: la escala total de datos es más crítica que una pequeña expansión del modelo}

La slide 62 compara por separado aproximadamente 1/20 de la escala de parámetros y 1/20 de la escala de datos. Los resultados enfatizan el rol de data scale: reducir datos degrada significativamente múltiples métricas, mientras que en este rango de modelos pequeños, modest model-size change tiene menor impacto. Aquí aún hay que evitar extrapolar a LLMs frontier; en escalas mayores, la proporción óptima entre parámetros, datos y compute podría ser distinta.

\lecturefigure{slide-062.jpg}{Contraste entre escala de datos y modelo: reducir datos causa degradación más evidente}{CS25 V6 official deck, p.~62}

\lecturefigure{slide-063.jpg}{Conclusión RQ1: no se detectó interferencia bilingüe universal}{CS25 V6 official deck, p.~63}

\lecturefigure{slide-064.jpg}{Conclusión RQ2: la estructura de speaker y code-switching tiene impacto menor al esperado}{CS25 V6 official deck, p.~64}

\lecturefigure{slide-065.jpg}{Conclusión RQ3: en el rango actual, data scale es más importante que modest model size}{CS25 V6 official deck, p.~65}

\lecturefigure{slide-066.jpg}{Takeaways de Bilingual BabyLM y límites experimentales}{CS25 V6 official deck, p.~66}

\teachervoice{00:25:25--00:29:01，el docente resume tres puntos: no hay interferencia monolingüe amplia; el impacto de la estructura de exposición es sorpresivamente débil; y la escala total de datos es más crítica que el tamaño del modelo en este rango. El docente simultáneamente limita las conclusiones a un small-scale setting.}

\begin{warningbox}{“El bilingüismo no interfiere” no es una ley condicional}
La distancia lingüística, el tokenizer, la proporción de mezcla, la duración del entrenamiento, las tareas generativas y la capacidad del modelo pueden cambiar los resultados. Lo que muestran los experimentos actuales es que bajo un conjunto controlado de condiciones small-scale no hay degeneración sistemática, no que todo multilingual training genere automáticamente transferencia positiva.
\end{warningbox}

Para la práctica de ingeniería, este conjunto de resultados se asemeja más a un orden de depuración que a una receta de mezcla. Si un modelo multilingüe degrada, primero revise el número efectivo de tokens por idioma, fertilidad del tokenizer, muestreo del batch y filtrado de evaluación (evaluation leakage), luego verifique competencia de capacidad. Speaker-split o estructura de code-switching pueden volverse importantes en tareas contextuales más fuertes; las métricas actuales de sintaxis y similitud lexical no son sensibles al cambio lingüístico a nivel de discurso. La prudencia del curso radica en separar "no se observó un efecto significativo" de "el efecto no existe", lo cual es el principio básico para leer cualquier resultado negativo.

\subsection{Resumen de la sección}

Ambos proyectos BabyLM demuestran conjuntamente que el valor de los experimentos small-scale reside en separar variables de datos. Baby Scale muestra que las propiedades de los datos familiares se correlacionan con el rendimiento de la tarea; Bilingual BabyLM indica que la mezcla lingüística no necesariamente causa interferencia y que la estructura de exposición y el volumen total de datos juegan roles distintos. Proporcionan pistas sobre mecanismos, no una fórmula universal para replicar directamente el aprendizaje infantil en modelos grandes.
